% Options for packages loaded elsewhere
\PassOptionsToPackage{unicode}{hyperref}
\PassOptionsToPackage{hyphens}{url}
\PassOptionsToPackage{dvipsnames,svgnames,x11names}{xcolor}
\documentclass[
  11pt,
  a4paper,
]{article}
\usepackage{xcolor}
\usepackage[margin=1in]{geometry}
\usepackage{amsmath,amssymb}
\setcounter{secnumdepth}{-\maxdimen} % remove section numbering
\usepackage{iftex}
\ifPDFTeX
  \usepackage[T1]{fontenc}
  \usepackage[utf8]{inputenc}
  \usepackage{textcomp} % provide euro and other symbols
\else % if luatex or xetex
  \usepackage{unicode-math} % this also loads fontspec
  \defaultfontfeatures{Scale=MatchLowercase}
  \defaultfontfeatures[\rmfamily]{Ligatures=TeX,Scale=1}
\fi
\usepackage{lmodern}
\ifPDFTeX\else
  % xetex/luatex font selection
  \ifLuaTeX
    \usepackage{luaotfload}
    \directlua{luaotfload.add_fallback("mainfontfallback",{
      "NotoSerifBengali:mode=harf"
    })}
  \fi
  \setmainfont[,RawFeature={fallback=mainfontfallback}]{Latin Modern Roman}
  \ifLuaTeX
    \usepackage{luaotfload}
    \directlua{luaotfload.add_fallback("sansfontfallback",{
      "NotoSansBengali:mode=harf"
    })}
  \fi
  \setsansfont[,RawFeature={fallback=sansfontfallback}]{Latin Modern Sans}
  \ifLuaTeX
    \usepackage{luaotfload}
    \directlua{luaotfload.add_fallback("monofontfallback",{
      "NotoSansBengali:mode=harf"
    })}
  \fi
  \setmonofont[,RawFeature={fallback=monofontfallback}]{DejaVu Sans Mono}
\fi
% Use upquote if available, for straight quotes in verbatim environments
\IfFileExists{upquote.sty}{\usepackage{upquote}}{}
\IfFileExists{microtype.sty}{% use microtype if available
  \usepackage[]{microtype}
  \UseMicrotypeSet[protrusion]{basicmath} % disable protrusion for tt fonts
}{}
\makeatletter
\@ifundefined{KOMAClassName}{% if non-KOMA class
  \IfFileExists{parskip.sty}{%
    \usepackage{parskip}
  }{% else
    \setlength{\parindent}{0pt}
    \setlength{\parskip}{6pt plus 2pt minus 1pt}}
}{% if KOMA class
  \KOMAoptions{parskip=half}}
\makeatother
\usepackage{color}
\usepackage{fancyvrb}
\newcommand{\VerbBar}{|}
\newcommand{\VERB}{\Verb[commandchars=\\\{\}]}
\DefineVerbatimEnvironment{Highlighting}{Verbatim}{commandchars=\\\{\}}
% Add ',fontsize=\small' for more characters per line
\usepackage{framed}
\definecolor{shadecolor}{RGB}{248,248,248}
\newenvironment{Shaded}{\begin{snugshade}}{\end{snugshade}}
\newcommand{\AlertTok}[1]{\textcolor[rgb]{0.94,0.16,0.16}{#1}}
\newcommand{\AnnotationTok}[1]{\textcolor[rgb]{0.56,0.35,0.01}{\textbf{\textit{#1}}}}
\newcommand{\AttributeTok}[1]{\textcolor[rgb]{0.13,0.29,0.53}{#1}}
\newcommand{\BaseNTok}[1]{\textcolor[rgb]{0.00,0.00,0.81}{#1}}
\newcommand{\BuiltInTok}[1]{#1}
\newcommand{\CharTok}[1]{\textcolor[rgb]{0.31,0.60,0.02}{#1}}
\newcommand{\CommentTok}[1]{\textcolor[rgb]{0.56,0.35,0.01}{\textit{#1}}}
\newcommand{\CommentVarTok}[1]{\textcolor[rgb]{0.56,0.35,0.01}{\textbf{\textit{#1}}}}
\newcommand{\ConstantTok}[1]{\textcolor[rgb]{0.56,0.35,0.01}{#1}}
\newcommand{\ControlFlowTok}[1]{\textcolor[rgb]{0.13,0.29,0.53}{\textbf{#1}}}
\newcommand{\DataTypeTok}[1]{\textcolor[rgb]{0.13,0.29,0.53}{#1}}
\newcommand{\DecValTok}[1]{\textcolor[rgb]{0.00,0.00,0.81}{#1}}
\newcommand{\DocumentationTok}[1]{\textcolor[rgb]{0.56,0.35,0.01}{\textbf{\textit{#1}}}}
\newcommand{\ErrorTok}[1]{\textcolor[rgb]{0.64,0.00,0.00}{\textbf{#1}}}
\newcommand{\ExtensionTok}[1]{#1}
\newcommand{\FloatTok}[1]{\textcolor[rgb]{0.00,0.00,0.81}{#1}}
\newcommand{\FunctionTok}[1]{\textcolor[rgb]{0.13,0.29,0.53}{\textbf{#1}}}
\newcommand{\ImportTok}[1]{#1}
\newcommand{\InformationTok}[1]{\textcolor[rgb]{0.56,0.35,0.01}{\textbf{\textit{#1}}}}
\newcommand{\KeywordTok}[1]{\textcolor[rgb]{0.13,0.29,0.53}{\textbf{#1}}}
\newcommand{\NormalTok}[1]{#1}
\newcommand{\OperatorTok}[1]{\textcolor[rgb]{0.81,0.36,0.00}{\textbf{#1}}}
\newcommand{\OtherTok}[1]{\textcolor[rgb]{0.56,0.35,0.01}{#1}}
\newcommand{\PreprocessorTok}[1]{\textcolor[rgb]{0.56,0.35,0.01}{\textit{#1}}}
\newcommand{\RegionMarkerTok}[1]{#1}
\newcommand{\SpecialCharTok}[1]{\textcolor[rgb]{0.81,0.36,0.00}{\textbf{#1}}}
\newcommand{\SpecialStringTok}[1]{\textcolor[rgb]{0.31,0.60,0.02}{#1}}
\newcommand{\StringTok}[1]{\textcolor[rgb]{0.31,0.60,0.02}{#1}}
\newcommand{\VariableTok}[1]{\textcolor[rgb]{0.00,0.00,0.00}{#1}}
\newcommand{\VerbatimStringTok}[1]{\textcolor[rgb]{0.31,0.60,0.02}{#1}}
\newcommand{\WarningTok}[1]{\textcolor[rgb]{0.56,0.35,0.01}{\textbf{\textit{#1}}}}
\usepackage{longtable,booktabs,array}
\usepackage{calc} % for calculating minipage widths
% Correct order of tables after \paragraph or \subparagraph
\usepackage{etoolbox}
\makeatletter
\patchcmd\longtable{\par}{\if@noskipsec\mbox{}\fi\par}{}{}
\makeatother
% Allow footnotes in longtable head/foot
\IfFileExists{footnotehyper.sty}{\usepackage{footnotehyper}}{\usepackage{footnote}}
\makesavenoteenv{longtable}
\setlength{\emergencystretch}{3em} % prevent overfull lines
\providecommand{\tightlist}{%
  \setlength{\itemsep}{0pt}\setlength{\parskip}{0pt}}
% paper_preamble.tex — shared LaTeX preamble for the PCMA/SI paper set.
%
% Compiled with lualatex (NOT pdflatex) because several papers contain
% Bengali script: the CLL paper realises one cognitive state into
% English, Bengali and Banglish, so the Bengali output is EVIDENCE and
% cannot be transliterated away without destroying the point being made.
%
% Font selection is NOT done here. pandoc >= 3 has native
% `mainfontfallback` / `monofontfallback` variables that emit the
% luaotfload fallback registration in the correct position relative to
% unicode-math; doing it by hand in this header lands after unicode-math
% and breaks font loading outright. See tools/build_papers.sh.

\usepackage{amsmath}
\usepackage{amssymb}
\usepackage{mathtools}   % \xleftrightarrow etc. used in cll.md §5
\usepackage{booktabs}

% ---------------------------------------------------------------------
% Unicode symbols appearing in PROSE rather than inside $...$.
%
% The papers were written as markdown and read as markdown, so they use
% real characters — "the ω-boost", "α-child", "ΔR coherence" — instead of
% LaTeX macros. Under unicode-math those codepoints are math-mode-only
% and abort the build ("\symbballowed only in math mode"). Mapping each
% to \ensuremath keeps the markdown readable on GitHub AND compilable,
% which is the whole point of keeping one source.
% ---------------------------------------------------------------------
\usepackage{newunicodechar}

% Greek, lower case
\newunicodechar{α}{\ensuremath{\alpha}}
\newunicodechar{β}{\ensuremath{\beta}}
\newunicodechar{γ}{\ensuremath{\gamma}}
\newunicodechar{δ}{\ensuremath{\delta}}
\newunicodechar{ε}{\ensuremath{\varepsilon}}
\newunicodechar{η}{\ensuremath{\eta}}
\newunicodechar{θ}{\ensuremath{\theta}}
\newunicodechar{κ}{\ensuremath{\kappa}}
\newunicodechar{λ}{\ensuremath{\lambda}}
\newunicodechar{μ}{\ensuremath{\mu}}
\newunicodechar{π}{\ensuremath{\pi}}
\newunicodechar{ρ}{\ensuremath{\rho}}
\newunicodechar{σ}{\ensuremath{\sigma}}
\newunicodechar{τ}{\ensuremath{\tau}}
\newunicodechar{φ}{\ensuremath{\varphi}}
\newunicodechar{χ}{\ensuremath{\chi}}
\newunicodechar{ψ}{\ensuremath{\psi}}
\newunicodechar{ω}{\ensuremath{\omega}}
\newunicodechar{ν}{\ensuremath{\nu}}
\newunicodechar{ξ}{\ensuremath{\xi}}
\newunicodechar{ζ}{\ensuremath{\zeta}}
\newunicodechar{ι}{\ensuremath{\iota}}

% Greek, upper case
\newunicodechar{Δ}{\ensuremath{\Delta}}
\newunicodechar{Φ}{\ensuremath{\Phi}}
\newunicodechar{Σ}{\ensuremath{\Sigma}}
\newunicodechar{Ψ}{\ensuremath{\Psi}}
\newunicodechar{Ω}{\ensuremath{\Omega}}
\newunicodechar{Γ}{\ensuremath{\Gamma}}
\newunicodechar{Λ}{\ensuremath{\Lambda}}
\newunicodechar{Θ}{\ensuremath{\Theta}}
\newunicodechar{Π}{\ensuremath{\Pi}}

% Operators and relations
\newunicodechar{×}{\ensuremath{\times}}
\newunicodechar{·}{\ensuremath{\cdot}}
\newunicodechar{±}{\ensuremath{\pm}}
\newunicodechar{−}{\ensuremath{-}}
\newunicodechar{≈}{\ensuremath{\approx}}
\newunicodechar{≤}{\ensuremath{\leq}}
\newunicodechar{≥}{\ensuremath{\geq}}
\newunicodechar{≠}{\ensuremath{\neq}}
\newunicodechar{∈}{\ensuremath{\in}}
\newunicodechar{∂}{\ensuremath{\partial}}
\newunicodechar{∇}{\ensuremath{\nabla}}
\newunicodechar{√}{\ensuremath{\sqrt{}}}
\newunicodechar{∞}{\ensuremath{\infty}}
\newunicodechar{∝}{\ensuremath{\propto}}
\newunicodechar{∑}{\ensuremath{\sum}}
\newunicodechar{∫}{\ensuremath{\int}}
\newunicodechar{ℝ}{\ensuremath{\mathbb{R}}}
\newunicodechar{ℕ}{\ensuremath{\mathbb{N}}}
\newunicodechar{‖}{\ensuremath{\|}}
\newunicodechar{⊕}{\ensuremath{\oplus}}
\newunicodechar{⊗}{\ensuremath{\otimes}}
\newunicodechar{∀}{\ensuremath{\forall}}
\newunicodechar{∃}{\ensuremath{\exists}}
\newunicodechar{¬}{\ensuremath{\neg}}
\newunicodechar{∧}{\ensuremath{\wedge}}
\newunicodechar{∨}{\ensuremath{\vee}}

% Arrows
\newunicodechar{→}{\ensuremath{\rightarrow}}
\newunicodechar{←}{\ensuremath{\leftarrow}}
\newunicodechar{↔}{\ensuremath{\leftrightarrow}}
\newunicodechar{↑}{\ensuremath{\uparrow}}
\newunicodechar{↓}{\ensuremath{\downarrow}}
\newunicodechar{⇒}{\ensuremath{\Rightarrow}}
\newunicodechar{⇐}{\ensuremath{\Leftarrow}}
\newunicodechar{⇔}{\ensuremath{\Leftrightarrow}}

% Super/subscripts written as literal characters
\newunicodechar{²}{\ensuremath{^{2}}}
\newunicodechar{³}{\ensuremath{^{3}}}
\newunicodechar{⁴}{\ensuremath{^{4}}}
\newunicodechar{⁵}{\ensuremath{^{5}}}
\newunicodechar{⁶}{\ensuremath{^{6}}}
\newunicodechar{₀}{\ensuremath{_{0}}}
\newunicodechar{₁}{\ensuremath{_{1}}}
\newunicodechar{₂}{\ensuremath{_{2}}}
\newunicodechar{₃}{\ensuremath{_{3}}}
\newunicodechar{ᵀ}{\ensuremath{^{\mathsf{T}}}}

% Status glyphs. The papers use emoji as status markers in tables
% (implemented / failed / warning). Mapped to typographic equivalents:
% a colour emoji font in lualatex is unreliable, and the MEANING is what
% carries — a reader needs to see "this one failed", not a red cross.
\newunicodechar{✅}{\ensuremath{\checkmark}}
\newunicodechar{✓}{\ensuremath{\checkmark}}
\newunicodechar{❌}{\ensuremath{\times}}
\newunicodechar{✗}{\ensuremath{\times}}
\newunicodechar{⚠}{\textbf{!}}
\newunicodechar{🧩}{\textbf{[piece]}}
\newunicodechar{📋}{\textbf{[list]}}
\newunicodechar{🟡}{\textbf{[partial]}}
\newunicodechar{🔴}{\textbf{[dead]}}
\newunicodechar{🟢}{\textbf{[live]}}
\newunicodechar{▼}{\ensuremath{\blacktriangledown}}
\newunicodechar{►}{\ensuremath{\blacktriangleright}}
\newunicodechar{□}{\ensuremath{\square}}
\newunicodechar{█}{\rule{0.6em}{0.9ex}}
\newunicodechar{•}{\textbullet}
\newunicodechar{†}{\textdagger}
\newunicodechar{‑}{-}
\newunicodechar{‐}{-}
\newunicodechar{…}{\ldots}
\newunicodechar{।}{\textbar}   % danda: rendered by the Bengali fallback in text
\newunicodechar{⁻}{\ensuremath{^{-}}}

% Zero-width formatting codepoints. These are invisible in an editor but
% are real characters to LaTeX: an emoji written as "⚠️" is TWO
% codepoints (U+26A0 U+FE0F), so declaring only the visible one leaves a
% variation selector with no definition and the build dies on a symbol
% nobody can see. Both are mapped to nothing.
\newunicodechar{^^^^fe0f}{}    % VARIATION SELECTOR-16
\newunicodechar{^^^^200d}{}    % ZERO WIDTH JOINER
\newunicodechar{^^^^200b}{}    % ZERO WIDTH SPACE

\usepackage{longtable}
\usepackage{array}

% Comparison matrices and cost ledgers are wide. Let them shrink to the
% text block rather than overflow the margin.
\usepackage{adjustbox}

% Code blocks: wrap long lines instead of running off the page. Several
% papers quote real C++ with long identifiers.
\usepackage{fvextra}
\fvset{breaklines=true,breakanywhere=true,fontsize=\small}

\setlength{\parskip}{0.5em}
\setlength{\parindent}{0pt}

\usepackage{fancyhdr}
\pagestyle{fancy}
\fancyhf{}
\fancyfoot[C]{\thepage}
\fancyhead[L]{\small\itshape Synthetic Intelligence — research lineage}
\fancyhead[R]{\small\itshape irfan@furylogic.com}
\renewcommand{\headrulewidth}{0.2pt}
\setlength{\headheight}{14pt}
\setlength{\headsep}{18pt}
\usepackage{bookmark}
\IfFileExists{xurl.sty}{\usepackage{xurl}}{} % add URL line breaks if available
\urlstyle{same}
\hypersetup{
  pdftitle={Bug Fix Roadmap: Four Phases to Complete SI},
  pdfauthor={Irfan Mahir --- Independent Researcher, Dhaka, Bangladesh --- irfan@furylogic.com},
  colorlinks=true,
  linkcolor={black},
  filecolor={Maroon},
  citecolor={Blue},
  urlcolor={[HTML]{0B5394}},
  pdfcreator={LaTeX via pandoc}}

\title{Bug Fix Roadmap: Four Phases to Complete SI}
\author{Irfan Mahir --- Independent Researcher, Dhaka, Bangladesh ---
irfan@furylogic.com}
\date{}

\begin{document}
\maketitle

{
\hypersetup{linkcolor=}
\setcounter{tocdepth}{3}
\tableofcontents
}
\section{Bug Fix Roadmap: Four Phases to Complete
SI}\label{bug-fix-roadmap-four-phases-to-complete-si}

\textbf{Irfan Mahir} \emph{Independent Researcher, Dhaka, Bangladesh}
irfan@furylogic.com

\begin{quote}
\emph{``The current architecture is solid enough to carry the entire
roadmap. No magic needed --- only phased engineering, cognition,
construction, and formalization.''}
\end{quote}

\textbf{Companion documents.} - Theory:
\texttt{research-paper/synthetic-intelligence.md} - Mechanism:
\texttt{research-paper/rbg.md} - Implementation:
\texttt{research-paper/failed\_but\_parts\_works/deprected/how\_we\_implement.md}
- Detailed per-item roadmap: \texttt{research-paper/future\_fix.md}

This document is the \emph{phase-organized} version of
\texttt{future\_fix.md}, grouping the ten remaining open problems into
four sequential phases by nature of the work (engineering → cognition →
construction → formalization).

\textbf{Every entry carries a ``⚠️ Caveat / Correction'' block} naming
the subtle implementation pitfall that would break the fix if ignored.
Read these before writing code. They were surfaced during a validation
pass after the initial roadmap was drafted.

Summary of caveats:

\begin{longtable}[]{@{}
  >{\raggedright\arraybackslash}p{(\linewidth - 2\tabcolsep) * \real{0.5000}}
  >{\raggedright\arraybackslash}p{(\linewidth - 2\tabcolsep) * \real{0.5000}}@{}}
\toprule\noalign{}
\begin{minipage}[b]{\linewidth}\raggedright
Item
\end{minipage} & \begin{minipage}[b]{\linewidth}\raggedright
Caveat
\end{minipage} \\
\midrule\noalign{}
\endhead
\bottomrule\noalign{}
\endlastfoot
1.1 Sparse tick & Wavefront saturates without a decay rule \\
1.2 Auto-threshold & Sign of \(k\) must be positive; use EMA to prevent
oscillation \\
1.3 Online L6-A & Must call \texttt{\_locked} variant to avoid mutex
deadlock \\
2.1 Syntax slots & Phase spacing is \(2\pi/N\), not \(\pi/N\) \\
2.2 Band-pass & Must be an added per-tunnel field; not emergent from
k/m \\
3.1 Multi-ontology & Use paired-cavities + auto-bridge, not shared
cavity \\
3.2 Reverse abduction & Physics is symmetric --- reverse is an
observation policy, not a mode \\
4.1 Liouville & State is a stratified space, not a smooth manifold \\
\end{longtable}

\begin{center}\rule{0.5\linewidth}{0.5pt}\end{center}

\subsection{Phase 1: Runtime Engineering (Days to
Weeks)}\label{phase-1-runtime-engineering-days-to-weeks}

\textbf{✅ STATUS: SHIPPED.} All three items (1.1, 1.2, 1.3) are
implemented, tested, and merged. See \texttt{data/phase1/results.md} for
the full report. 20-seed benchmark drift: 0.07\% RMS (well under ±5\%
budget). Snapshot: \texttt{.snapshots/pre\_phase1/}.

These bugs can be fixed directly in the current C++ codebase with
existing infrastructure. No new theory required.

\subsubsection{1.1 Sparse Structural Tick (10⁴--10⁵
concepts)}\label{sparse-structural-tick-10ux207410ux2075-concepts}

\textbf{Method: Active Wavefront Tracking.}

The engine currently scans all N cavities every physics tick ---
\(O(N)\) per iteration. This is fine at 200 concepts but fails at
10⁴--10⁵.

\textbf{Fix.} Maintain a dynamic \texttt{touched\_nodes} set inside
\texttt{PhysicsEngine}. Only cavities whose displacement changed by more
than \(\varepsilon\) in the last tick enter the next tick. The wavefront
propagates naturally: touched cavities activate their tunnels, which
touch neighbouring cavities, which enter the working set for the next
tick.

\begin{Shaded}
\begin{Highlighting}[]
\BuiltInTok{std::}\NormalTok{unordered\_set}\OperatorTok{\textless{}}\DataTypeTok{int}\OperatorTok{\textgreater{}}\NormalTok{ touched\_cavities}\OperatorTok{;}
\DataTypeTok{float}\NormalTok{ wavefront\_eps }\OperatorTok{=} \FloatTok{1e{-}4}\BuiltInTok{f}\OperatorTok{;}

\DataTypeTok{void}\NormalTok{ tick}\OperatorTok{()} \OperatorTok{\{}
    \CommentTok{// Only iterate over touched cavities and their incident tunnels.}
    \CommentTok{// At tick end: any cavity with |Δx| \textgreater{} wavefront\_eps enters}
    \CommentTok{// touched\_cavities for the next tick; others age out after K ticks.}
\OperatorTok{\}}
\end{Highlighting}
\end{Shaded}

\textbf{Files:} \texttt{src/physics.hpp} \textbf{Success test:}
Bit-identical results on the 200-concept ontology; sub-5-ms tick time on
a synthetic 10⁵-concept graph.

\textbf{⚠️ Caveat --- wavefront saturation under sustained driving.}
Under a full training run (\texttt{train\_multi\ metformin,berberine\ 5}
drives at amplitude 8 for 40 ticks), the wavefront can propagate to
\textasciitilde40\% of the ontology within a single round. If cavities
are added to \texttt{touched\_nodes} but never removed, the working set
only grows and sparse tick degrades back to the dense O(N) case.

\textbf{Required companion mechanism --- decay rule.} A cavity leaves
\texttt{touched\_nodes} after \texttt{K\_idle} ticks with no fresh input
(a fresh input = tunnel force above \texttt{wavefront\_eps} OR external
drive). Without this, the mechanism is not truly sparse --- it is only
lazily initialised. With decay it is genuinely
\(O(|\text{active wavefront}|)\).

\begin{Shaded}
\begin{Highlighting}[]
\KeywordTok{struct}\NormalTok{ WavefrontEntry }\OperatorTok{\{}
    \DataTypeTok{int}\NormalTok{ cavity\_id}\OperatorTok{;}
    \DataTypeTok{int}\NormalTok{ quiet\_ticks}\OperatorTok{;}   \CommentTok{// reset to 0 on fresh input; increment each tick}
\OperatorTok{\};}
\AttributeTok{const} \DataTypeTok{int}\NormalTok{ K\_idle }\OperatorTok{=} \DecValTok{20}\OperatorTok{;}   \CommentTok{// remove from touched set after 20 quiet ticks}
\end{Highlighting}
\end{Shaded}

\begin{center}\rule{0.5\linewidth}{0.5pt}\end{center}

\subsubsection{1.2 Auto-Tuning Fission
Thresholds}\label{auto-tuning-fission-thresholds}

\textbf{Method: Homeostatic Energy Feedback.}

Currently \texttt{fission\_I\_thresh\ =\ 0.15} is hardcoded. On a new
ontology it may fire cascade or not fire at all.

\textbf{Fix.} Replace the constant with a feedback function of system
energy:

\[
I_\text{thresh}(t) = I_0 + k \cdot E_\text{system}(t)
\]

Where \(E_\text{system}\) is the current total mechanical energy from
\texttt{physics.audit\_energy()}. If fission events cause a sudden
energy spike, the threshold auto-raises and blocks spurious follow-up
fissions. If the system is quiet, the threshold drops and allows the
mechanism to fire on genuine contradictions.

\begin{Shaded}
\begin{Highlighting}[]
\DataTypeTok{void}\NormalTok{ refresh\_fission\_thresholds}\OperatorTok{()} \OperatorTok{\{}
    \DataTypeTok{float}\NormalTok{ E }\OperatorTok{=}\NormalTok{ physics}\OperatorTok{.}\NormalTok{audit\_energy}\OperatorTok{().}\NormalTok{total\_mechanical}\OperatorTok{;}
\NormalTok{    fission\_I\_thresh }\OperatorTok{=} \FloatTok{0.10}\BuiltInTok{f} \OperatorTok{+} \FloatTok{0.005}\BuiltInTok{f} \OperatorTok{*}\NormalTok{ E}\OperatorTok{;}
\OperatorTok{\}}
\end{Highlighting}
\end{Shaded}

\textbf{Files:} \texttt{src/layer2.hpp} \textbf{Success test:} Fission
count stable per training round across ontologies of varying size; no
cascade on null controls; no missed contradictions on active controls.

\textbf{⚠️ Caveat --- sign discipline.} \texttt{k} must be strictly
\emph{positive}. Negative feedback (rising E → rising threshold → harder
to fission) suppresses cascade. A sign flip produces positive feedback
(rising E → falling threshold → easier to fission → cascade collapse of
the entire graph within one training round). This is not a subtle bug
--- a wrong sign kills the run in \textasciitilde50 ticks. Guard with a
compile-time
\texttt{static\_assert(k\_energy\_feedback\ \textgreater{}\ 0.0f)} or a
runtime sanity check.

\textbf{⚠️ Caveat --- response timescale.} If \(E_\text{system}\) is
instantaneous, one fission spikes \(E\), the threshold rises, next tick
\(E\) falls, the threshold drops, another fission fires --- oscillation.
Use an EMA (matches the existing \texttt{sil.tension\_integral} idiom in
\texttt{rifa.hpp}):

\[
\bar{E}(t) = 0.95 \cdot \bar{E}(t-1) + 0.05 \cdot E(t)
\]

\begin{Shaded}
\begin{Highlighting}[]
\DataTypeTok{float}\NormalTok{ E\_smooth }\OperatorTok{=} \FloatTok{0.0}\BuiltInTok{f}\OperatorTok{;}
\DataTypeTok{void}\NormalTok{ refresh\_fission\_thresholds}\OperatorTok{()} \OperatorTok{\{}
    \DataTypeTok{float}\NormalTok{ E }\OperatorTok{=}\NormalTok{ physics}\OperatorTok{.}\NormalTok{audit\_energy}\OperatorTok{().}\NormalTok{total\_mechanical}\OperatorTok{;}
\NormalTok{    E\_smooth }\OperatorTok{=} \FloatTok{0.95}\BuiltInTok{f} \OperatorTok{*}\NormalTok{ E\_smooth }\OperatorTok{+} \FloatTok{0.05}\BuiltInTok{f} \OperatorTok{*}\NormalTok{ E}\OperatorTok{;}
\NormalTok{    fission\_I\_thresh }\OperatorTok{=} \FloatTok{0.10}\BuiltInTok{f} \OperatorTok{+} \FloatTok{0.005}\BuiltInTok{f} \OperatorTok{*}\NormalTok{ E\_smooth}\OperatorTok{;}   \CommentTok{// k \textgreater{} 0 required}
\OperatorTok{\}}
\end{Highlighting}
\end{Shaded}

Without smoothing the threshold flips between fires on adjacent ticks
and the streak counter never reaches
\texttt{fission\_hold\_target\ =\ 3}.

\begin{center}\rule{0.5\linewidth}{0.5pt}\end{center}

\subsubsection{1.3 Online L6-A Discovery}\label{online-l6-a-discovery}

\textbf{Method: In-Loop Intersection.}

\texttt{discover\_analogies()} is currently invoked from the shell as a
separate command. In a live-running system it should fire the moment two
fissioned α-children accumulate ≥ K shared targets.

\textbf{Fix.} Call \texttt{discover\_analogies\_locked()} inside the
structural tick, right after a fission event fires. Only run
intersection for the newly-fissioned α against every existing α (not the
full N² pairs).

\begin{Shaded}
\begin{Highlighting}[]
\ControlFlowTok{if} \OperatorTok{(}\NormalTok{fire\_id }\OperatorTok{\textgreater{}=} \DecValTok{0}\OperatorTok{)} \OperatorTok{\{}
\NormalTok{    sub}\OperatorTok{{-}\textgreater{}}\NormalTok{fission\_locked}\OperatorTok{(}\NormalTok{fire\_id}\OperatorTok{,} \OperatorTok{...);}
    \DataTypeTok{int}\NormalTok{ new\_alpha }\OperatorTok{=}\NormalTok{ sub}\OperatorTok{{-}\textgreater{}}\NormalTok{fission\_map}\OperatorTok{[}\NormalTok{fire\_id}\OperatorTok{].}\NormalTok{first}\OperatorTok{;}
    \CommentTok{// Check this new alpha against every existing alpha.}
\NormalTok{    check\_analogies\_for}\OperatorTok{(}\NormalTok{new\_alpha}\OperatorTok{);}
\OperatorTok{\}}
\end{Highlighting}
\end{Shaded}

\textbf{Files:} \texttt{src/layer2.hpp} \textbf{Success test:} Bridges
appear within one structural tick of the second concept fissioning;
total bridge count matches the post-hoc \texttt{discover\_analogies}
invocation.

\textbf{⚠️ Caveat --- mutex discipline.} Inside
\texttt{structural\_tick\_once()} the physics mutex
\texttt{sub-\textgreater{}mtx} is already held. Calling the public
\texttt{discover\_analogies()} there would try to re-acquire the same
mutex → deadlock. Instead, split the method into a public
\texttt{discover\_analogies()} (acquires mutex) and an internal
\texttt{discover\_analogies\_locked()} (assumes caller holds the mutex),
mirroring the existing \texttt{fission()} / \texttt{fission\_locked()}
split. Only the \texttt{\_locked} variant is called from within the
structural tick.

\begin{Shaded}
\begin{Highlighting}[]
\CommentTok{// Public — safe from the shell.}
\DataTypeTok{int}\NormalTok{ discover\_analogies}\OperatorTok{(}\DataTypeTok{int}\NormalTok{ min\_shared }\OperatorTok{=} \DecValTok{2}\OperatorTok{,} \OperatorTok{...)} \OperatorTok{\{}
    \BuiltInTok{std::}\NormalTok{lock\_guard}\OperatorTok{\textless{}}\BuiltInTok{std::}\NormalTok{mutex}\OperatorTok{\textgreater{}}\NormalTok{ lk}\OperatorTok{(}\NormalTok{sub}\OperatorTok{{-}\textgreater{}}\NormalTok{mtx}\OperatorTok{);}
    \ControlFlowTok{return}\NormalTok{ discover\_analogies\_locked}\OperatorTok{(}\NormalTok{min\_shared}\OperatorTok{,} \OperatorTok{...);}
\OperatorTok{\}}
\CommentTok{// Internal — assumes caller holds sub{-}\textgreater{}mtx.}
\DataTypeTok{int}\NormalTok{ discover\_analogies\_locked}\OperatorTok{(}\DataTypeTok{int}\NormalTok{ min\_shared}\OperatorTok{,} \OperatorTok{...);}
\end{Highlighting}
\end{Shaded}

\begin{center}\rule{0.5\linewidth}{0.5pt}\end{center}

\subsection{Phase 2: Cognitive Interface
(Weeks)}\label{phase-2-cognitive-interface-weeks}

\textbf{✅ STATUS: SHIPPED.} Both items (2.1 syntax macro-layer, 2.2
band-pass per-tunnel property) implemented and merged. See
\texttt{data/phase2/results.md}. 20-seed benchmark bit-identical to
Phase 1. Snapshot: \texttt{.snapshots/pre\_phase2/}.

This phase connects the engine to human language and gives L6-A a
principled frequency law.

\subsubsection{2.1 Language Stage 2 --- Syntax
Macro-Layer}\label{language-stage-2-syntax-macro-layer}

\textbf{Method: Kuramoto Syntax Hubs.}

Add a second macro-layer to \texttt{physics.hpp} with fixed
grammatical-role frequencies --- one hub per slot (Subject, Verb,
Object). When the semantic wave from Layer 1 reaches these hubs,
phase-locking (from HNCA §2) determines which concept lands in which
slot.

\textbf{Fix.} New struct + one extra Kuramoto pass + a temporal readout
head that samples the hubs in canonical order (θ\_S = 0, θ\_V = 2π/3,
θ\_O = 4π/3):

\begin{Shaded}
\begin{Highlighting}[]
\KeywordTok{struct}\NormalTok{ SyntaxSlot }\OperatorTok{\{}
    \DataTypeTok{float}\NormalTok{ natural\_omega}\OperatorTok{;}   \CommentTok{// fixed per slot}
    \DataTypeTok{float}\NormalTok{ phase}\OperatorTok{,}\NormalTok{ amplitude}\OperatorTok{;}
    \BuiltInTok{std::}\NormalTok{string role}\OperatorTok{;}
    \DataTypeTok{int}\NormalTok{ locked\_cavity\_id}\OperatorTok{;}
\OperatorTok{\};}
\BuiltInTok{std::}\NormalTok{vector}\OperatorTok{\textless{}}\NormalTok{SyntaxSlot}\OperatorTok{\textgreater{}}\NormalTok{ syntax\_slots}\OperatorTok{;}

\BuiltInTok{std::}\NormalTok{string physics\_generate\_sentence}\OperatorTok{(}\DataTypeTok{int}\NormalTok{ probe\_cavity\_id}\OperatorTok{)} \OperatorTok{\{}
    \CommentTok{// Drive, run, sample slots at canonical phases,}
    \CommentTok{// concatenate labels in Subject{-}Verb{-}Object order.}
\OperatorTok{\}}
\end{Highlighting}
\end{Shaded}

\textbf{Files:} \texttt{src/physics.hpp}, \texttt{src/mal.hpp},
\texttt{src/nl\_translator.hpp} \textbf{Success test:}
\texttt{generate\ metformin} produces \texttt{metformin\ inhibits\ mtor}
--- deterministic under seed, correct edge direction, no next-token
prediction anywhere in the pipeline.

\textbf{⚠️ Correction --- phase spacing must be \(2\pi/N\), not
\(\pi/N\).} For an \(N\)-slot syntax layer the slot phases must be
evenly distributed around the unit circle at
\(\{0, 2\pi/N, 4\pi/N, \ldots, 2\pi(N-1)/N\}\). Using
\(\{0, \pi/3, 2\pi/3\}\) instead of \(\{0, 2\pi/3, 4\pi/3\}\) compresses
the Subject and Verb slots into the same half-plane; the readout head
cannot distinguish them because their phase-locking windows overlap.

\begin{itemize}
\tightlist
\item
  \textbf{3 slots (SVO):} \(\theta_S = 0\), \(\theta_V = 2\pi/3\),
  \(\theta_O = 4\pi/3\)
\item
  \textbf{4 slots (SVO+Modifier):} \(\{0, \pi/2, \pi, 3\pi/2\}\)
\item
  \textbf{5 slots (SVO+Modifier+Adverb):}
  \(\{0, 2\pi/5, 4\pi/5, 6\pi/5, 8\pi/5\}\)
\item
  \textbf{General N:} \(\theta_k = 2\pi k / N\) for \(k \in [0, N)\)
\end{itemize}

Pick the natural frequency \(\omega\) per slot as harmonic ratios of a
base frequency \(\omega_0\) (e.g.~\(\omega_0, 2\omega_0, 3\omega_0\)) to
keep the Kuramoto lock structure clean and orthogonal.

\begin{center}\rule{0.5\linewidth}{0.5pt}\end{center}

\subsubsection{2.2 Rubber-Tunnel Frequency Selectivity
Theorem}\label{rubber-tunnel-frequency-selectivity-theorem}

\textbf{Method: Band-Pass Transfer Functions.}

The empirical 3 Hz probe frequency for L6-A was picked by inspection.
Rubber tunnels should not transmit all frequencies indiscriminately ---
only a specific band should resonate.

\textbf{Fix.} Derive the transfer function of a mass-spring-damper
tunnel:

\[
H(\omega) = \frac{k}{\sqrt{(k - m\omega^2)^2 + (c\omega)^2}}
\]

Compute the optimal probe frequency algorithmically from a tunnel's
material parameters:

\begin{Shaded}
\begin{Highlighting}[]
\DataTypeTok{float}\NormalTok{ optimal\_probe\_freq}\OperatorTok{(}\AttributeTok{const}\NormalTok{ Tunnel}\OperatorTok{\&}\NormalTok{ t}\OperatorTok{,} \DataTypeTok{float}\NormalTok{ m }\OperatorTok{=} \FloatTok{1.0}\BuiltInTok{f}\OperatorTok{)} \OperatorTok{\{}
    \ControlFlowTok{return} \BuiltInTok{std::}\NormalTok{sqrt}\OperatorTok{(}\NormalTok{t}\OperatorTok{.}\NormalTok{stiffness }\OperatorTok{/}\NormalTok{ m}\OperatorTok{)} \OperatorTok{/} \OperatorTok{(}\FloatTok{2.0}\BuiltInTok{f} \OperatorTok{*}\NormalTok{ M\_PI}\OperatorTok{);}
\OperatorTok{\}}
\end{Highlighting}
\end{Shaded}

Then \texttt{repurpose\ \textless{}label\textgreater{}\ auto} inspects
the target's incident tunnel materials and picks the frequency that
maximizes the ratio \texttt{H\_rubber(ω)\ /\ H\_steel(ω)} --- the
analogy-signal-to-mechanistic-noise ratio.

\textbf{Files:} \texttt{src/physics.hpp}, new
\texttt{research-paper/mims\_transfer.md} \textbf{Success test:}
Reproduce the atorvastatin-at-rank-10 result without manually specifying
3 Hz.

\textbf{⚠️ Correction --- band-pass must be an added per-tunnel
property, not emergent from rubber's k/m.} The naive expectation is that
a rubber tunnel (\(k=0.05, m=1\)) naturally resonates at
\(\omega_0 = \sqrt{k/m} \approx 0.22\) rad/s ≈ 0.036 Hz --- nowhere near
the empirical 3 Hz used in the L6-A demo. The rubber material as
currently specified is \emph{not} a band-pass filter at 3 Hz; it is a
low-pass filter with corner far below 3 Hz.

\textbf{Fix.} Add a per-tunnel resonance frequency as a separate field,
not derived from \(k\) and \(m\):

\begin{Shaded}
\begin{Highlighting}[]
\KeywordTok{struct}\NormalTok{ Tunnel }\OperatorTok{\{}
    \CommentTok{// existing fields (stiffness, damping, polarity, base\_stiffness,}
    \CommentTok{//                  plasticity\_delta, cavity\_a, cavity\_b) ...}
    \DataTypeTok{float}\NormalTok{ resonance\_omega }\OperatorTok{=} \FloatTok{0.0}\BuiltInTok{f}\OperatorTok{;}   \CommentTok{// 0 = pass{-}through; nonzero = band{-}pass at ω}
    \DataTypeTok{float}\NormalTok{ Q\_factor        }\OperatorTok{=} \FloatTok{3.0}\BuiltInTok{f}\OperatorTok{;}   \CommentTok{// resonance sharpness (dimensionless)}
\OperatorTok{\};}
\end{Highlighting}
\end{Shaded}

Then the tunnel's force contribution gets a frequency-selective gate:

\begin{Shaded}
\begin{Highlighting}[]
\DataTypeTok{float}\NormalTok{ bp\_gate }\OperatorTok{=} \OperatorTok{(}\NormalTok{t}\OperatorTok{.}\NormalTok{resonance\_omega }\OperatorTok{\textgreater{}} \FloatTok{0.0}\BuiltInTok{f}\OperatorTok{)}
              \OperatorTok{?}\NormalTok{ Q\_factor\_gain}\OperatorTok{(}\NormalTok{sim\_time}\OperatorTok{,}\NormalTok{ t}\OperatorTok{.}\NormalTok{resonance\_omega}\OperatorTok{,}\NormalTok{ t}\OperatorTok{.}\NormalTok{Q\_factor}\OperatorTok{)}
              \OperatorTok{:} \FloatTok{1.0}\BuiltInTok{f}\OperatorTok{;}
\NormalTok{F\_A }\OperatorTok{*=}\NormalTok{ bp\_gate}\OperatorTok{;}
\NormalTok{F\_B }\OperatorTok{*=}\NormalTok{ bp\_gate}\OperatorTok{;}
\end{Highlighting}
\end{Shaded}

where \texttt{Q\_factor\_gain} is a standard band-pass response centred
on \texttt{resonance\_omega}. When
\texttt{discover\_analogies\_locked()} inserts a rubber bridge, it sets
\texttt{resonance\_omega\ =\ 3\ Hz} explicitly. Steel-material tunnels
leave \texttt{resonance\_omega\ =\ 0} and pass all frequencies.

This makes band-pass filtering an explicit \emph{L6-A design choice},
not an implicit consequence of a material parameter that does not
actually produce it.

\begin{center}\rule{0.5\linewidth}{0.5pt}\end{center}

\subsection{Phase 3: Cross-Domain Construction
(Months)}\label{phase-3-cross-domain-construction-months}

\textbf{✅ STATUS: SHIPPED.} Both items (3.1 multi-ontology +
cross-domain L6-A, 3.2 abduction discovery) implemented and merged.
Bonus: dataset-driven axiom extraction
(\texttt{tools/extract\_axioms.py}) --- 462 real axioms harvested from
C4 + Pile via Hearst patterns. See \texttt{data/phase3/results.md}.
Engine grew from 216→1020 nodes. Snapshot:
\texttt{.snapshots/pre\_phase3\_4/}.

This phase teaches the engine to reason across independent ontologies
and discover new inference rule families.

\subsubsection{3.1 Cross-Domain Transfer (Syllogism →
Algebra)}\label{cross-domain-transfer-syllogism-algebra}

\textbf{Method: Multi-Ontology Substrates.}

Currently the engine loads one ontology at a time. To demonstrate
transfer, load two --- say \texttt{data/syllogism.axioms} and
\texttt{data/algebra.axioms} --- and share the physics substrate.
Concepts that appear in both (e.g.~\texttt{implies}, \texttt{equals})
become bridge cavities; L6-A discovers cross-domain rubber-bridges
between structurally analogous concepts (\texttt{increase}
\textasciitilde{} \texttt{addition}, \texttt{associativity}
\textasciitilde{} \texttt{transitive\_composition}).

\textbf{Fix.} - Add \texttt{domain\_id} field to \texttt{Cavity} (0 =
shared foundation, 1 = syllogism, 2 = algebra). - Train the substrate on
syllogism traces, measure baseline algebra recall, then train jointly
and measure improvement. - Extend L6-A discovery to prefer cross-domain
bridges (higher structural information gain per bridge).

\textbf{Files:} \texttt{src/physics.hpp}, new
\texttt{data/algebra.axioms}, extended \texttt{src/rade.hpp}
\textbf{Success test:} Recall of algebra concepts improves after
syllogism training via shared logical connectives; L6-A surfaces at
least one cross-domain analogy
(\texttt{associativity\ \textasciitilde{}\ transitive\_composition}).

\textbf{⚠️ Design decision --- shared cavity vs paired cavities with
bridge.} When a concept name appears in both files (say
\texttt{implies}), you must choose one of two behaviours:

\begin{itemize}
\tightlist
\item
  \textbf{Option A --- Shared cavity.} Unified label lookup:
  \texttt{syllogism.implies} and \texttt{algebra.implies} become the
  \emph{same} cavity. Simpler. Standard graph-database practice.
  Cross-domain transfer is trivial (both domains drive the same cavity
  directly).
\item
  \textbf{Option B --- Paired cavities + auto-bridge.} Each ontology
  gets its own \texttt{implies} cavity (\texttt{syllogism.implies},
  \texttt{algebra.implies}). L6-A discovers the analogy at runtime and
  inserts a rubber bridge. Type distinction is preserved.
\end{itemize}

\textbf{Recommended: Option B} for the paper. It makes the cross-domain
L6-A claim non-trivial --- the bridge is \emph{found}, not asserted. It
also matches Axiom I of the SI paper (typed graph rewriting:
\texttt{syllogism.implies} and \texttt{algebra.implies} have distinct
types until L6-A relates them).

Implementation:

\begin{Shaded}
\begin{Highlighting}[]
\KeywordTok{struct}\NormalTok{ Cavity }\OperatorTok{\{}
    \CommentTok{// existing fields ...}
    \DataTypeTok{int}\NormalTok{ domain\_id }\OperatorTok{=} \DecValTok{0}\OperatorTok{;}   \CommentTok{// 0 = shared foundation, 1 = syllogism, 2 = algebra}
\OperatorTok{\};}

\CommentTok{// Axiom loader keys labels by (domain\_id, label) rather than label alone.}
\BuiltInTok{std::}\NormalTok{string domain\_prefix}\OperatorTok{(}\DataTypeTok{int}\NormalTok{ d}\OperatorTok{)} \OperatorTok{\{}
    \ControlFlowTok{return} \OperatorTok{(}\NormalTok{d }\OperatorTok{==} \DecValTok{0}\OperatorTok{)} \OperatorTok{?} \StringTok{""} \OperatorTok{:} \OperatorTok{(}\NormalTok{d }\OperatorTok{==} \DecValTok{1} \OperatorTok{?} \StringTok{"syllogism."} \OperatorTok{:} \StringTok{"algebra."}\OperatorTok{);}
\OperatorTok{\}}
\end{Highlighting}
\end{Shaded}

L6-A discovery then naturally surfaces
\texttt{syllogism.implies\ \textasciitilde{}\ algebra.implies} as its
first cross-domain bridge --- a testable prediction the paper can
report.

\begin{center}\rule{0.5\linewidth}{0.5pt}\end{center}

\subsubsection{3.2 Abduction Rule Discovery from
Primitives-Only}\label{abduction-rule-discovery-from-primitives-only}

\textbf{Method: Reverse Wave Propagation.}

Currently abduction is a hand-authored primitive. To \emph{discover}
abduction rules, drive the physics engine backwards: given a goal cavity
\(G\), propagate a wave from \(G\) \emph{backwards} along its incoming
\texttt{implies} tunnels. If the wave reaches a primitive \(P\) with
sufficient amplitude, save the path \(P \Rightarrow G\) as a new
abduction rule.

\textbf{Fix.} - Add a
\texttt{drive\_reverse(node\_id,\ freq,\ amp,\ dur)} call that injects
energy against the polarity direction of incident tunnels. - After
running the reverse simulation, walk the amplitude field back from
\(G\); every path with cumulative amplitude above threshold becomes a
candidate abduction rule. - Verify each candidate via the existing
Layer-2 verifier (does forward evaluation of \(P \Rightarrow G\)
actually derive \(G\)?).

\textbf{Files:} new \texttt{src/abduction\_discovery.hpp}, extended
\texttt{src/physics.hpp} \textbf{Success test:} Given
\texttt{raining\ implies\ wet\_ground} in the ontology and only
\texttt{wet\_ground} as a driven observation, the system spawns
\texttt{hypothesis(raining)} as a candidate abduction that verifies
cleanly.

\textbf{⚠️ Correction --- the physics is symmetric; ``reverse'' is not a
sign flip.} The tunnel force equation is \(F_{A \to B} = -F_{B \to A}\)
by Newton's third law. Driving cavity \(A\) automatically
back-propagates into cavity \(B\) if a tunnel connects them, because the
force is transmitted both ways. There is no separate ``reverse mode'' to
switch on. Attempting to ``drive backwards'' by flipping the sign of
\texttt{F\_A} would violate momentum conservation and destabilize the
physics.

\textbf{Correct implementation --- reverse is an observation policy on
forward physics.} Drive the goal cavity at normal amplitude. Let the
physics settle normally (all forces are symmetric --- the wave
\emph{already} back-propagates along incoming tunnels). Then observe
which \emph{upstream} cavities (cavities with excitatory tunnels ending
at the goal) accumulate the most residual energy after settling ---
those are the hypothesis candidates.

\begin{Shaded}
\begin{Highlighting}[]
\DataTypeTok{void}\NormalTok{ discover\_abductions}\OperatorTok{(}\DataTypeTok{int}\NormalTok{ goal\_id}\OperatorTok{,} \DataTypeTok{float}\NormalTok{ threshold }\OperatorTok{=} \FloatTok{0.02}\BuiltInTok{f}\OperatorTok{)} \OperatorTok{\{}
\NormalTok{    physics}\OperatorTok{.}\NormalTok{reset\_field}\OperatorTok{();}
\NormalTok{    physics}\OperatorTok{.}\NormalTok{drive}\OperatorTok{(}\NormalTok{goal\_id}\OperatorTok{,} \FloatTok{10.0}\BuiltInTok{f}\OperatorTok{,} \FloatTok{6.0}\BuiltInTok{f}\OperatorTok{,} \DecValTok{20}\OperatorTok{);}
\NormalTok{    physics}\OperatorTok{.}\NormalTok{run}\OperatorTok{(}\DecValTok{200}\OperatorTok{);}

    \CommentTok{// Scan cavities with incoming excitatory tunnels ending at goal\_id.}
    \ControlFlowTok{for} \OperatorTok{(}\AttributeTok{const} \KeywordTok{auto}\OperatorTok{\&}\NormalTok{ t }\OperatorTok{:}\NormalTok{ physics}\OperatorTok{.}\NormalTok{tunnels}\OperatorTok{)} \OperatorTok{\{}
        \ControlFlowTok{if} \OperatorTok{(}\NormalTok{t}\OperatorTok{.}\NormalTok{polarity }\OperatorTok{\textless{}} \FloatTok{0.0}\BuiltInTok{f}\OperatorTok{)} \ControlFlowTok{continue}\OperatorTok{;}                  \CommentTok{// excitatory only}
        \ControlFlowTok{if} \OperatorTok{(}\NormalTok{t}\OperatorTok{.}\NormalTok{cavity\_b }\OperatorTok{!=}\NormalTok{ goal\_id}\OperatorTok{)} \ControlFlowTok{continue}\OperatorTok{;}              \CommentTok{// ends at goal}
        \DataTypeTok{int}\NormalTok{ p }\OperatorTok{=}\NormalTok{ t}\OperatorTok{.}\NormalTok{cavity\_a}\OperatorTok{;}                                \CommentTok{// upstream candidate}
        \DataTypeTok{float}\NormalTok{ amp\_p }\OperatorTok{=} \BuiltInTok{std::}\NormalTok{fabs}\OperatorTok{(}\NormalTok{physics}\OperatorTok{.}\NormalTok{cavities}\OperatorTok{[}\NormalTok{p}\OperatorTok{].}\NormalTok{displacement}\OperatorTok{);}
        \ControlFlowTok{if} \OperatorTok{(}\NormalTok{amp\_p }\OperatorTok{\textgreater{}}\NormalTok{ threshold}\OperatorTok{)} \OperatorTok{\{}
\NormalTok{            register\_abduction\_candidate}\OperatorTok{(}\NormalTok{p}\OperatorTok{,}\NormalTok{ goal\_id}\OperatorTok{);}
            \CommentTok{// Verify via Layer{-}2: does forward evaluation of p ⇒ goal derive goal?}
            \ControlFlowTok{if} \OperatorTok{(}\NormalTok{verifier}\OperatorTok{.}\NormalTok{walk\_back}\OperatorTok{(}\NormalTok{p}\OperatorTok{,}\NormalTok{ goal\_id}\OperatorTok{))} \OperatorTok{\{}
\NormalTok{                save\_as\_new\_rule}\OperatorTok{(}\NormalTok{p}\OperatorTok{,}\NormalTok{ goal\_id}\OperatorTok{);}
            \OperatorTok{\}}
        \OperatorTok{\}}
    \OperatorTok{\}}
\OperatorTok{\}}
\end{Highlighting}
\end{Shaded}

Save this observation as a new rule via the existing meta-learner +
verifier. No new physics primitive is needed --- reverse abduction is a
\emph{policy} on forward dynamics, not a separate mode.

\begin{center}\rule{0.5\linewidth}{0.5pt}\end{center}

\subsection{Phase 4: Mathematical Formalization (Months to
Years)}\label{phase-4-mathematical-formalization-months-to-years}

\textbf{✅ STATUS: SHIPPED.} All three theorem papers drafted: -
\texttt{research-paper/liouville\_analog.md} --- divergence-free proof -
\texttt{research-paper/graph\_fisher.md} --- Boltzmann Fisher metric
(computable from \texttt{plasticity\_delta}) + L6-A optimality theorem +
shell command \texttt{fisher\_distance} -
\texttt{research-paper/cr\_bound.md} ---
\texttt{D\_max\ ≤\ exp(β·T\_max·P)} capstone

Fisher metric is empirically checkable via
\texttt{fisher\_distance\ \textless{}a\textgreater{}\ \textless{}b\textgreater{}}
in the shell. See \texttt{data/phase4/results.md}.

This phase is theorems, not code. It is what turns the empirical work
into a paper Nature / Science can review.

\subsubsection{4.1 Formal Liouville Analog for Tension-Bounded
Expansion}\label{formal-liouville-analog-for-tension-bounded-expansion}

\textbf{Method: Phase-Space Analysis.}

Classical Liouville: phase-space volume is preserved under Hamiltonian
flow. RB+G expands the phase space under contradiction pressure --- the
classical theorem does not apply.

\textbf{Conjecture.} Define an extended phase space
\(\mathcal{P} = \bigsqcup_N (\mathbb{R}^{2N} \times \mathcal{G}_N)\)
where \(\mathcal{G}_N\) is the space of typed graphs on \(N\) vertices.
The weighted volume

\[
V(t) = \sum_N \int_{\mathcal{G}_N} e^{-\beta T(G)} \, d\text{vol}
\]

is preserved by the combined Hamiltonian flow + fission dynamics, where
\(T(G)\) is the RIFA-tension functional and \(\beta\) is a ``cognitive
temperature''.

\textbf{Proof sketch (to be completed).} - Piecewise-Hamiltonian flow is
measure-preserving on each stratum (standard). - Fission jump preserves
weighted volume across strata --- the exponential weight
\(e^{-\beta \Delta T}\) exactly compensates for the added 3 degrees of
freedom (2 cavities + 1 tunnel).

\textbf{Files:} new \texttt{research-paper/liouville\_analog.md}
\textbf{Success test:} A theorem statement + rigorous proof + empirical
corroboration (measure \(V(t)\) across the 20-seed benchmark; should
stay within 5\%).

\textbf{⚠️ Caveat --- the state is a stratified space, not a naive
smooth manifold.} The claim ``prove the graph-valued state is a
continuous differentiable manifold'' understates the technical work.
Under fission, the phase space gains 3 dimensions discontinuously (2 new
cavities + 1 new tunnel). The combined space is not smooth in the
classical differential-geometric sense --- it is a \textbf{stratified
space} (each \(N\)-stratum is a smooth manifold; fission jumps are
controlled discontinuities between strata).

Two established categorical approaches to handle this:

\begin{itemize}
\tightlist
\item
  \textbf{Nerve construction.} Embed the typed graph category into a
  simplicial complex, then take its geometric realisation. The result is
  a smooth manifold with singular strata at each fission event. Standard
  Grothendieck construction machinery. Well-defined but produces a large
  ambient space.
\item
  \textbf{Configuration-space embedding.} Identify each graph \(G\) with
  the tuple \((N(G), \text{adjacency}(G), \text{labels}(G))\) and treat
  each \(N\)-stratum \(\mathcal{G}_N\) as a separate smooth manifold,
  with the fission jump as a controlled discontinuity between
  \(\mathcal{G}_N\) and \(\mathcal{G}_{N+2}\). Requires stratified Morse
  theory to handle the transitions rigorously.
\end{itemize}

Both approaches are legitimate; neither is a one-line proof.
\textbf{Realistic timeframe:} 6--12 months of graduate-level
differential geometry + category theory work. This is where the roadmap
crosses from ``engineering'' into ``genuine research programme''.

\textbf{✅ Solution Sketch (Mahir).} An independent solution path via
\textbf{divergence-free Hamiltonian flow} --- cleaner than the
stratified-space approach above, and reaches the same weighted-volume
conservation theorem.

Define the Hamiltonian directly over the substrate state:

\[
H(x_i, v_i, \theta_k, \omega_k) = \sum_i \tfrac{1}{2} m_i v_i^2
                                + \tfrac{1}{2} \sum_{(i,j) \in E} k_{ij}(x_i - p_{ij} x_j)^2
                                + \text{(HNCA phase terms)}
\]

The engine's canonical pairs are already in the code: \((x_i, v_i)\) per
cavity and \((\theta_k, \omega_k)\) per macro-node. These are conjugate
positions and momenta.

\textbf{Step 1 --- Between fission events (divergence-free flow).} Show
that the velocity field satisfies

\[
\nabla \cdot \vec{v} = \sum_i \left(\frac{\partial \dot{x}_i}{\partial x_i} + \frac{\partial \dot{v}_i}{\partial v_i}\right) = 0
\]

by Hamilton's equations \(\dot{x}_i = \partial H/\partial v_i\),
\(\dot{v}_i = -\partial H/\partial x_i\), and equality of mixed
partials. This is the classical proof --- one page --- and establishes
that the \(N\)-stratum volume is conserved between fissions.

\textbf{Step 2 --- At fission (weighted-volume bridging).} At the
fission event, the state gains 3 dimensions (2 cavities + 1 tunnel). The
tension functional \(T(G)\) drops by \(\Delta T\) (the contradiction is
resolved into the α/β pair). Choose the cognitive temperature \(\beta\)
such that the exponential weight \(e^{-\beta \Delta T}\) exactly cancels
the added volume:

\[
\text{added volume} = \Delta V_\text{new dims} = c_0 \cdot \text{(vol of 3-D fission cell)},
\] \[
\beta = \frac{1}{\Delta T} \ln\!\left(\frac{V_\text{after}}{V_\text{before}}\right).
\]

The weighted volume \(V(t) = \int e^{-\beta T} d\text{vol}\) is
preserved across the jump by construction.

\textbf{Combined theorem.} Under Hamiltonian flow + tension-gated
fission, \(V(t) = \int_{\mathcal{P}} e^{-\beta T} \, d\text{vol}\) is
invariant.

This is a \emph{simpler} proof than the stratified Morse-theory approach
and avoids the categorical machinery of nerve constructions. It uses
only Hamilton's equations + one exponential weight. Publishable as a
stand-alone theorem in classical Hamiltonian mechanics language.

\textbf{Effort under this approach:} 3--4 months of physics-flavoured
derivation, not 12 months of category theory.

\begin{center}\rule{0.5\linewidth}{0.5pt}\end{center}

\subsubsection{4.2 Fisher-Information Metric on Graph-Valued
States}\label{fisher-information-metric-on-graph-valued-states}

\textbf{Method: Continuous Differentiable Manifold Embedding.}

The Fisher-information metric is defined for parametric families
\(\{p_\theta\}\) on smooth manifolds. For graph-valued states (Axiom I
of SI) we need a probability distribution over graphs.

\textbf{Approach.} Define an exponential family over typed graphs:

\[
p_\theta(G) = \frac{1}{Z(\theta)} \exp\left(-\beta \sum_v \phi(x_v, \theta) - \gamma \sum_e \psi(x_{v(e)}, x_{w(e)}, \theta)\right)
\]

parametrised by the concept-cavity displacements. Prove this has a
well-defined Fisher metric on the parameter space, and show that the
RB+G expansion trajectory is a Fisher-geodesic minimizing the tension
integral \(\int T(G_t) dt\).

\textbf{Files:} new \texttt{research-paper/graph\_fisher.md}
\textbf{Success test:} Theorem statement + proof of geodesic optimality
+ categorical functoriality under typed-graph morphisms.

\textbf{✅ Solution Sketch (Mahir).} Use the \textbf{standard Fisher
metric} on a concrete Boltzmann distribution computable from quantities
the engine already reports.

\textbf{Step 1 --- Pick a distribution.} The physics engine tracks total
mechanical energy \(E(G, \theta)\) via \texttt{physics.audit\_energy()}.
Define the Boltzmann distribution over graphs:

\[
p(G | \theta) = \frac{1}{Z(\theta)} \exp\left(-\beta \, E(G, \theta)\right),
\]

parametrised by cavity displacements and tunnel plasticities collected
into \(\theta\).

\textbf{Step 2 --- Compute the Fisher metric.} For any exponential
family, the Fisher metric has the closed form:

\[
g_{\mu\nu}(\theta) = \mathbb{E}\!\left[\frac{\partial \log p}{\partial \theta^\mu} \cdot \frac{\partial \log p}{\partial \theta^\nu}\right] = \beta^2 \, \mathrm{Cov}\!\left(\frac{\partial E}{\partial \theta^\mu}, \frac{\partial E}{\partial \theta^\nu}\right).
\]

The right-hand side is a \emph{covariance of energy derivatives} --- and
the engine already tracks this via \texttt{corr\_ema} (tunnel pair
correlation) in \texttt{layer2.hpp}. \textbf{The Fisher metric is
directly computable from the existing telemetry.} No new data structure
needed.

\begin{Shaded}
\begin{Highlighting}[]
\CommentTok{// After training, extract Fisher metric from corr\_ema:}
\NormalTok{Matrix fisher\_metric}\OperatorTok{()} \OperatorTok{\{}
    \DataTypeTok{int}\NormalTok{ T }\OperatorTok{=} \OperatorTok{(}\DataTypeTok{int}\OperatorTok{)}\NormalTok{sub}\OperatorTok{{-}\textgreater{}}\NormalTok{tunnels}\OperatorTok{.}\NormalTok{size}\OperatorTok{();}
\NormalTok{    Matrix g}\OperatorTok{(}\NormalTok{T}\OperatorTok{,}\NormalTok{ T}\OperatorTok{);}
    \DataTypeTok{float}\NormalTok{ beta\_sq }\OperatorTok{=} \FloatTok{1.0}\BuiltInTok{f}\OperatorTok{;}   \CommentTok{// pick β once, tune empirically}
    \ControlFlowTok{for} \OperatorTok{(}\DataTypeTok{int}\NormalTok{ i }\OperatorTok{=} \DecValTok{0}\OperatorTok{;}\NormalTok{ i }\OperatorTok{\textless{}}\NormalTok{ T}\OperatorTok{;} \OperatorTok{++}\NormalTok{i}\OperatorTok{)}
        \ControlFlowTok{for} \OperatorTok{(}\DataTypeTok{int}\NormalTok{ j }\OperatorTok{=} \DecValTok{0}\OperatorTok{;}\NormalTok{ j }\OperatorTok{\textless{}}\NormalTok{ T}\OperatorTok{;} \OperatorTok{++}\NormalTok{j}\OperatorTok{)}
\NormalTok{            g}\OperatorTok{(}\NormalTok{i}\OperatorTok{,}\NormalTok{ j}\OperatorTok{)} \OperatorTok{=}\NormalTok{ beta\_sq }\OperatorTok{*}\NormalTok{ corr\_ema}\OperatorTok{[}\NormalTok{i}\OperatorTok{]} \OperatorTok{*}\NormalTok{ corr\_ema}\OperatorTok{[}\NormalTok{j}\OperatorTok{];}
    \ControlFlowTok{return}\NormalTok{ g}\OperatorTok{;}
\OperatorTok{\}}
\end{Highlighting}
\end{Shaded}

\textbf{Step 3 --- Analogy-bridge optimality claim.} L6-A inserts a
rubber bridge between α-children \(a_i, a_j\). The bridge is
\emph{optimal in the Fisher sense} iff the Fisher distance
\(d_g(a_i, a_j)\) is small --- i.e. the two cavities are statistically
indistinguishable under \(p(G|\theta)\). This gives a hard mathematical
criterion for when L6-A should fire:

\[
\text{insert bridge}(a_i, a_j) \iff d_g(a_i, a_j) < \varepsilon_\text{Fisher}.
\]

\textbf{Publishable claim:} \emph{``L6-A analogical bridging is a
Fisher-optimal compression of the graph state --- bridges are inserted
exactly when the Fisher metric declares the endpoints statistically
equivalent.''}

This is a much stronger claim than ``L6-A finds real analogies in the
insulin ontology'' --- it says L6-A is \emph{the} correct mechanism
under an information-geometric criterion.

\textbf{Effort:} 3--4 months, largely writing + one C++ helper function
to extract the metric from \texttt{corr\_ema}.

\begin{center}\rule{0.5\linewidth}{0.5pt}\end{center}

\subsubsection{4.3 Upper Bound on Thinker Index
C/R}\label{upper-bound-on-thinker-index-cr}

\textbf{Method:} Depends on 4.1 and 4.2.

Given the weighted phase-space volume (4.1) and the Fisher metric (4.2),
the reachable derivation count within \(L\) composition steps is bounded
by the volume of a Fisher-ball of radius \(L \cdot \sqrt{T_\text{avg}}\)
around the seed graph.

\textbf{Conjecture.} For any physically realizable SI system with
bounded tension budget \(T_\text{max}\):

\[
C/R \leq \exp(\beta \cdot T_\text{max} \cdot |\Pi|)
\]

where \(|\Pi|\) is the primitive count and \(\beta\) is the cognitive
temperature from 4.1.

\textbf{Interpretation.} SI is not Turing-omnipotent within its budget.
The physical cost of derivation caps the depth of the reachable
derivation tree. This makes SI \emph{bounded rational} --- precisely
like a human brain constrained by metabolic energy.

\textbf{Files:} capstone theorem in the same paper as 4.1 and 4.2.
\textbf{Success test:} A finite upper bound formula + empirical
verification (as more primitives are added, does measured C/R stay under
the bound?).

\textbf{✅ Solution Sketch (Mahir).} Prove the exponential topological
bound directly from the branching structure of the derivation tree.

\textbf{Step 1 --- Notation cleanup.} The SI paper §3.4 uses
\(C(\mathcal{S})\) = composition rate, \(R(\mathcal{S})\) = retrieval
rate. To avoid clash, introduce two new symbols specifically for this
bound:

\begin{itemize}
\tightlist
\item
  \(P = |\Pi|\) --- number of stored axioms / primitives.
\item
  \(D\) --- number of reachable derivations (compositions of length
  \(\leq L\)).
\end{itemize}

Then \(C/R\) (thinker index of §3.4) is a function of \(D/P\): an SI
system with \(D \gg P\) has high \(C/R\).

\textbf{Step 2 --- Branching count.} At each fission event, the graph
gains 2 new cavities and 1 new tunnel --- each of which can participate
in future derivations. The branching factor \(k\) per fission satisfies
\(k \leq P\) (bounded by the primitive count). After \(L\) fissions, the
reachable derivation set grows as \(D \leq k^L\).

\textbf{Step 3 --- Tension budget cap.} Each fission event costs tension
\(\Delta T\). Under the RB+G generation cap \(g_\text{max} = 1\), the
number of fission events per training round is bounded by
\(L_\text{max} = T_\text{max} / \Delta T\). Substituting:

\[
D \leq k^{L_\text{max}} \leq P^{T_\text{max} / \Delta T}.
\]

Taking logarithms and using
\(\ln P \cdot T_\text{max} / \Delta T = \beta T_\text{max} P'\) with
\(P' = \ln P\):

\[
\boxed{D_\text{max} \leq \exp\!\left(\beta \cdot T_\text{max} \cdot P\right)}
\]

\textbf{Step 4 --- Interpretation.} This bound is \emph{loose but
finite} --- it says any SI system with bounded tension budget
\(T_\text{max}\) has finite reachable derivation count. It distinguishes
SI from Turing-omnipotent computation: derivation depth is physically
capped by tension budget, not by axiom count alone.

\textbf{Human comparison.} A biological brain has \(P \approx 10^{11}\)
concepts and metabolic tension budget \(T_\text{max} \approx 20\)W
(glucose supply). The bound predicts a finite reachable derivation set
--- bounded rationality, not unlimited cognition. This matches empirical
observation.

\textbf{Publishable statement:} \emph{``For any physically realizable SI
system, \(D_\text{max} \leq \exp(\beta \, T_\text{max} \, |\Pi|)\). SI
is bounded rational within its metabolic budget, not
Turing-omnipotent.''}

\textbf{Effort:} 2--3 months once 4.1 and 4.2 are drafted (this is the
capstone theorem --- a page-long combinatorial argument built on the
Liouville and Fisher results).

\begin{center}\rule{0.5\linewidth}{0.5pt}\end{center}

\subsection{Priority Ordering}\label{priority-ordering}

If executed sequentially:

\begin{longtable}[]{@{}
  >{\raggedright\arraybackslash}p{(\linewidth - 6\tabcolsep) * \real{0.2500}}
  >{\raggedright\arraybackslash}p{(\linewidth - 6\tabcolsep) * \real{0.2500}}
  >{\raggedright\arraybackslash}p{(\linewidth - 6\tabcolsep) * \real{0.2500}}
  >{\raggedright\arraybackslash}p{(\linewidth - 6\tabcolsep) * \real{0.2500}}@{}}
\toprule\noalign{}
\begin{minipage}[b]{\linewidth}\raggedright
Phase
\end{minipage} & \begin{minipage}[b]{\linewidth}\raggedright
Items
\end{minipage} & \begin{minipage}[b]{\linewidth}\raggedright
Time
\end{minipage} & \begin{minipage}[b]{\linewidth}\raggedright
Deliverable
\end{minipage} \\
\midrule\noalign{}
\endhead
\bottomrule\noalign{}
\endlastfoot
\textbf{Phase 1} & 1.1, 1.2, 1.3 & \textasciitilde2 weeks & Self-tuning,
self-scaling, self-discovering engine \\
\textbf{Phase 2} & 2.1, 2.2 & \textasciitilde1 month & Physics-driven
language generation with principled frequency law \\
\textbf{Phase 3} & 3.1, 3.2 & \textasciitilde6 months & Cross-domain
reasoning + abduction rule discovery \\
\textbf{Phase 4} & 4.1, 4.2, 4.3 & \textbf{\textasciitilde9--10 months}
(was \textasciitilde2 years --- updated after Mahir's divergence-free +
Boltzmann sketches) & Theorem-strength foundations \\
\end{longtable}

Phase 1 alone gives a demonstrably self-contained engine that scales to
10⁵ concepts. Phase 1 + Phase 2 gives an engine that generates language.
Phase 1 + Phase 2 + Phase 3 is complete SI \emph{in the engineering
sense}. Phase 4 is what makes it publishable at Nature-level.

\subsubsection{Parallelizable}\label{parallelizable}

\begin{itemize}
\tightlist
\item
  Phase 1 items 1.1, 1.2, 1.3 can be done concurrently.
\item
  Phase 2 items 2.1 and 2.2 are independent.
\item
  Phase 3 items 3.1 and 3.2 are independent.
\item
  Phase 4 items 4.1 and 4.2 are independent research threads; 4.3
  depends on both.
\end{itemize}

\begin{center}\rule{0.5\linewidth}{0.5pt}\end{center}

\subsection{What This Roadmap
Confirms}\label{what-this-roadmap-confirms}

\begin{itemize}
\tightlist
\item
  The current architecture is \textbf{solid enough} to carry the entire
  roadmap. No structural rewrite required.
\item
  No item requires abandoning the tri-layer architecture, the physics
  substrate, or the storage-floor honesty of the Translation Capacity
  Lemma.
\item
  Every item is either engineering, empirical research, or mathematical
  proof --- none are impossibility proofs.
\item
  Phase 1 alone can be completed in a fortnight and materially improves
  the engine's autonomy.
\item
  Phase 4 is where the framework goes from ``novel implementation'' to
  ``publishable theorem'' --- but the empirical work (Phases 1--3) does
  not depend on Phase 4 being finished first.
\end{itemize}

\begin{center}\rule{0.5\linewidth}{0.5pt}\end{center}

\textbf{Bug Fix Roadmap: four phases, no magic. Runtime engineering,
then cognitive interface, then cross-domain construction, then
mathematical formalization. The architecture holds. What remains is time
and work.}

\end{document}
